%
% Figures for the research report
% Included via \input{figures} from main.tex
%

% ============================================================
% Figure 1: Conceptual Framework
% ============================================================
\begin{figure}[t]
\centering
\begin{tikzpicture}[
  node distance=1.4cm and 0.8cm,
  box/.style={rectangle, rounded corners=3pt, draw=#1!60, fill=#1!12, text=#1!80!black, minimum width=2.6cm, minimum height=0.75cm, align=center, font=\footnotesize\sffamily, line width=0.5pt},
  widebox/.style={rectangle, rounded corners=3pt, draw=#1!60, fill=#1!12, text=#1!80!black, minimum width=8.5cm, minimum height=0.75cm, align=center, font=\footnotesize\sffamily\bfseries, line width=0.5pt},
  smallbox/.style={rectangle, rounded corners=2pt, draw=#1!60, fill=#1!10, text=#1!80!black, minimum width=2.1cm, minimum height=0.6cm, align=center, font=\footnotesize\sffamily, line width=0.4pt},
  arrow/.style={->, >=stealth, line width=0.6pt, draw=black!55},
  darrow/.style={<->, >=stealth, line width=0.7pt, draw=black!55},
]
  % Layer 1: Protocols
  \node[widebox=blue, label={[font=\footnotesize\sffamily\bfseries]north:Data Sources}] (proto) {Aave (V2/V3) \hspace{1cm} Compound (V2/V3) \hspace{1cm} MakerDAO};

  % Layer 2: Behavioral variables  
  \node[widebox=teal, below=0.8cm of proto, label={[font=\footnotesize\sffamily\bfseries]north:Behavioral Process Extraction}] (behave) {\textbf{Borrower Active Adjustment} \quad (collateral changes, debt management, response timing)};

  % Layer 3: Two tracks
  \node[box=violet, below left=1.3cm and -1.5cm of behave, minimum width=3.5cm] (rq1) {\textbf{RQ1: Behavioral Layer}\\Do borrowers systematically adjust\\as HF $\rightarrow$ 1.0?};
  \node[box=violet, below right=1.3cm and -1.5cm of behave, minimum width=3.5cm] (rq2) {\textbf{RQ2: Credit Layer}\\Does adjustment behavior predict\\future liquidations?};

  \node[box=red!80!black, below=0.7cm of rq1, minimum width=3.5cm, font=\footnotesize\sffamily] (out1) {Prospect Theory Validation\\{\small(loss aversion, reference-point}};
  \node[box=red!80!black, below=0.7cm of rq2, minimum width=3.5cm, font=\footnotesize\sffamily] (out2) {Behavioral Early Warning\\{\small(credit risk signals)}};

  % Arrows
  \draw[arrow, thick] (proto.south) -- (behave.north);
  \draw[arrow, thick] (behave.south) |- (rq1.north);
  \draw[arrow, thick] (behave.south) |- (rq2.north);
  \draw[arrow] (rq1.south) -- (out1.north);
  \draw[arrow] (rq2.south) -- (out2.north);
  \draw[darrow, thick] (rq1.east) -- node[above, font=\footnotesize\sffamily]{H2a: Incremental predictive power?} (rq2.west);

  % Left labels
  \node[font=\footnotesize\sffamily, text=black!60, rotate=90, above left=0.0cm and -1.8cm of behave, anchor=south] {$\longleftarrow$ Behavioral Finance $\longrightarrow$};

\end{tikzpicture}
\caption{Conceptual framework for studying borrower behavior and credit signals in DeFi lending. The framework proceeds from raw on-chain data (top) through behavioral process variable extraction (center) to the research questions (bottom). RQ1 characterizes how borrowers adjust positions as risk accumulates; RQ2 tests whether these behavioral patterns provide incremental predictive power for future liquidation events beyond conventional risk indicators.}
\label{fig:framework}
\end{figure}


% ============================================================
% Figure 2: Methodology Pipeline
% ============================================================
\begin{figure}[t]
\centering
\begin{tikzpicture}[
  node distance=0.9cm and 0.6cm,
  stage/.style={rectangle, rounded corners=4pt, draw=#1!70, fill=#1!15, text=black!80, minimum width=2.8cm, minimum height=0.8cm, align=center, font=\footnotesize\sffamily, line width=0.5pt},
  substage/.style={rectangle, rounded corners=2pt, draw=#1!50, fill=#1!8, text=black!70, minimum width=2.4cm, minimum height=0.55cm, align=center, font=\footnotesize\sffamily, line width=0.4pt},
  arrow/.style={->, >=stealth, line width=0.7pt, draw=black!60},
  label/.style={font=\footnotesize\sffamily\bfseries},
]
  % Stage 1
  \node[stage=blue] (s1) {1. Data Extraction\\{\small Ethereum RPC / Dune Analytics}};
  
  % Stage 2
  \node[stage=cyan, right=0.4cm of s1] (s2) {2. Event Classification\\{\small Active vs.\ Passive}};
  
  % Stage 3  
  \node[stage=teal, right=0.4cm of s2] (s3) {3. Trajectory\\Reconstruction\\{\small Daily HF, bands}};
  
  % Stage 4
  \node[stage=violet, right=0.4cm of s3] (s4) {4. Variable\\Construction\\{\small BH process vars}};
  
  % Stage 5
  \node[stage=red!70!black, right=0.4cm of s4] (s5) {5. Model\\Estimation\\{\small XGBoost + Logit}};

  % Arrows between stages
  \draw[arrow, thick] (s1.east) -- (s2.west);
  \draw[arrow, thick] (s2.east) -- (s3.west);
  \draw[arrow, thick] (s3.east) -- (s4.west);
  \draw[arrow, thick] (s4.east) -- (s5.west);

  % Substages under stages
  \node[substage=blue, below=0.2cm of s1] (s1a) {Borrow, Repay,\\Deposit, Withdraw};
  \node[substage=cyan, below=0.2cm of s2] (s2a) {msg.sender\\classification};
  \node[substage=teal, below=0.2cm of s3] (s3a) {min HF, band\\days, drawdown};
  \node[substage=violet, below=0.2cm of s4] (s4a) {adjust rate,\\latency, intensity};
  \node[substage=red!70!black, below=0.2cm of s5] (s5a) {AUC-ROC, PR-AUC\\DeLong test};

  % Output arrow
  \node[font=\small\sffamily, below=0.6cm of s5a, text=black!70] {$\longrightarrow$ Liquidation Prediction (7d / 30d / 90d)};

\end{tikzpicture}
\caption{Methodology pipeline from raw blockchain data to predictive model evaluation. Stage 1 extracts transaction events from Ethereum; Stage 2 classifies transactions as active (borrower-initiated) or passive (liquidation bot); Stage 3 reconstructs position risk trajectories; Stage 4 constructs behavioral process variables; and Stage 5 estimates and compares predictive models.}
\label{fig:methodology}
\end{figure}


% ============================================================
% Figure 3: Risk Trajectory / Health Factor Illustration
% ============================================================
\begin{figure}[t]
\centering
\begin{tikzpicture}[
  scale=1.0,
  >=stealth,
  axisline/.style={->, line width=0.6pt},
  dataline/.style={line width=0.8pt},
]
  % Axes
  \draw[axisline] (0,0) -- (10.5,0) node[right, font=\footnotesize] {Time $\longrightarrow$};
  \draw[axisline] (0,0) -- (0,5.5) node[above, font=\footnotesize] {Health Factor};
  
  % Labels
  \node[font=\footnotesize] at (5.0,-0.5) {Observation Window (daily frequency)};
  
  % Y-axis labels
  \draw[dashed, gray!40] (0,0.0) -- (9.0,0.0);
  \node[font=\footnotesize, left] at (-0.15,0.0) {0};
  \draw[dashed, gray!40] (0,1.0) -- (9.0,1.0);
  \node[font=\footnotesize, left] at (0,1.0) {1.0};
  \node[font=\footnotesize, red!70!black] at (-0.9, 0.55) {\shortstack{Liquidation\\Threshold}};
  
  \draw[dashed, gray!40] (0,2.0) -- (9.0,2.0);
  \node[font=\footnotesize, left] at (0,2.0) {2.0};
  \draw[dashed, gray!40] (0,3.0) -- (9.0,3.0);
  \node[font=\footnotesize, left] at (0,3.0) {3.0};
  \draw[dashed, gray!40] (0,4.0) -- (9.0,4.0);
  \node[font=\footnotesize, left] at (0,4.0) {4.0};
  
  % Risk bands (shaded)
  \fill[red!8] (0,1.0) rectangle (9.0,0.0);
  \fill[orange!8] (0,2.0) rectangle (9.0,1.0);
  \fill[yellow!8] (0,3.0) rectangle (9.0,2.0);
  \fill[green!8] (0,5.0) rectangle (9.0,3.0);
  
  \node[font=\footnotesize\sffamily, text=red!50!black, opacity=0.6] at (1.5,0.5) {Liquidation Zone};
  \node[font=\footnotesize\sffamily, text=orange!50!black, opacity=0.6] at (1.5,1.5) {High Risk (HF $<$ 2)};
  \node[font=\footnotesize\sffamily, text=yellow!50!black, opacity=0.6] at (1.5,2.5) {Moderate Risk};
  \node[font=\footnotesize\sffamily, text=green!50!black, opacity=0.6] at (1.5,4.0) {Safe Zone (HF $>$ 3)};
  
  % Simulated HF trajectory
  \draw[dataline, blue!60, line width=1.1pt] plot[smooth] coordinates {
    (0,3.5) (0.4,3.3) (0.8,3.1) (1.2,2.8) (1.6,2.5) (2.0,2.3) 
    (2.4,2.0) (2.6,1.7) (2.7,1.5) (2.85,1.25) 
  };
  
  % Active adjustment events (stars)
  \node[star, star points=5, star point ratio=2.25, draw=green!70!black, fill=green!60, minimum size=2mm, inner sep=0pt] at (1.6,2.5) {};
  \draw[->, green!60!black, line width=0.5pt] (1.6,2.9) -- (1.6,2.5) node[pos=0, above, font=\tiny] {+collateral};
  
  \node[star, star points=5, star point ratio=2.25, draw=green!70!black, fill=green!60, minimum size=2mm, inner sep=0pt] at (2.4,2.0) {};
  \draw[->, green!60!black, line width=0.5pt] (2.4,2.4) -- (2.4,2.0) node[pos=0, above, font=\tiny] {repay debt};
  
  % Liquidation event (cross)
  \node[circle, draw=red!80!black, fill=red!60, minimum size=2.5mm, inner sep=0pt] at (2.85,1.25) {\tiny\color{white}\textbf{!}};
  \draw[->, red!80!black, line width=0.5pt] (3.2,1.25) -- (2.95,1.25) node[pos=0, right, font=\tiny] {liquidation};
  
  % Legend
  \node[font=\footnotesize] at (6.5,4.5) {
    \shortstack[l]{
      \tikz{\draw[blue!60, line width=1.1pt] (0,0)--(0.5,0);} \,\,\, HF trajectory\\
      \tikz{\node[star, star points=5, star point ratio=2.25, fill=green!60, minimum size=1.5mm, inner sep=0pt] at (0,0) {};} \,\,\, Active adjustment\\
      \tikz{\node[circle, fill=red!60, minimum size=1.5mm, inner sep=0pt] at (0,0) {};} \,\,\, Liquidation event
    }
  };

\end{tikzpicture}
\caption{Illustrative health factor (HF) trajectory and behavioral response. A borrowing position's HF declines toward the liquidation threshold (HF = 1.0) as collateral value decreases. Active borrower adjustments (green stars) represent collateral additions or debt repayment. The central empirical question is whether the pattern of these adjustments—their timing, magnitude, and frequency—provides predictive information about eventual liquidation beyond what is captured by the HF itself.}
\label{fig:trajectory}
\end{figure}
